\documentclass[11pt,a4paper]{article}

\usepackage[margin=2.5cm]{geometry}
\usepackage{hyperref}
\usepackage{amsmath,amssymb}
\usepackage{parskip}
\usepackage{microtype}
\usepackage{titlesec}
\usepackage{booktabs}
\usepackage{enumitem}

\titleformat{\section}{\large\bfseries}{\thesection}{0.6em}{}
\titleformat{\subsection}{\normalsize\bfseries}{\thesubsection}{0.6em}{}

\newcommand{\status}[1]{\hfill{\small\textsc{#1}}}

\title{\textbf{Experiment plan: open items}}
\author{}
\date{2026-11-01}

\begin{document}
\maketitle

This is the list of changes the experiment plan in \texttt{paper.tex}
needs before Stage~3 starts, from a review of \texttt{paper.tex} and
\texttt{methodology.tex} on 2026-11-01. Each item names the problem, where
it shows, and the fix. Items are ordered by how much they change what gets
run. The ones in \S\ref{sec:run} change the runs; the ones in
\S\ref{sec:docs} change only the documents. Every item is \textsc{open}
until ticked off here.

Terms used below. An \emph{arm} is one configuration in the comparison
(A to F in the paper's ablation table). A \emph{leaf} is one scored
trajectory. \emph{Fast success} means the final run succeeds within
$1.25\times$ the time-optimal batch time. The \emph{probe rate} is the
fraction of first replies that are probe code rather than a controller.
The \emph{blind reference} is the best hand-written controller that does
not know the hidden limit; the \emph{oracle} is the one that does.

\section{What holds up}

Not to be changed.
\begin{itemize}[leftmargin=1.5em]
  \item The gating logic. Stage~0 proves the hidden limit is worth knowing
  before any model runs; Stage~0b proves the runner carries the solution;
  Stages~1 and~2 set the floor and the training target; Stage~3 checks the
  plumbing against a known ceiling. Nothing downstream rests on an
  unverified step.
  \item The ablation isolates the claim. E against D tests the rollout
  tree; F against E tests the per-turn baseline; the trained arms match on
  leaves, and both equal-rollout and equal-wall-clock comparisons are
  reported.
  \item Fast success as the primary metric. Raw success is saturated by any
  slow cautious controller.
  \item The environment. The limit leaves no trace until it is crossed, so a
  probe is the only way to learn it, and Stage~0 showed the gap between
  blind and oracle is wide.
  \item The immediate next step. Re-run Stages~1 and~2 under the optional
  probe with the batch record and 1000-token turns, then decide on
  Stage~3. Items~\ref{it:wilcoxon} to~\ref{it:checkpoint} should be settled
  before Stage~3 starts, since they change what is run and how many seeds
  it needs.
\end{itemize}

\section{Items that change what gets run}
\label{sec:run}

\begin{enumerate}[leftmargin=1.8em, itemsep=0.8em]

\item \textbf{The statistical test cannot reach significance at three
seeds.} \status{open}
\label{it:wilcoxon}
\\
\emph{Problem.} The evaluation section specifies a paired Wilcoxon
signed-rank test over per-seed fast-success rates with at least three
seeds per arm. The test ranks the per-seed differences between two arms;
with $n$ pairs the smallest possible two-sided $p$ is $2/2^{n}$, so three
seeds cannot get below $0.25$.
\begin{center}
\small
\begin{tabular}{@{}cc@{}}
\toprule
Seeds per arm & Smallest two-sided $p$ \\
\midrule
3 & 0.25 \\
4 & 0.125 \\
5 & 0.0625 \\
6 & 0.031 \\
\bottomrule
\end{tabular}
\end{center}
\emph{Fix.} One of: (a) six seeds per arm for the arms in the headline
comparison (D and F at least); (b) pair over seed-by-batch units, so the
test has seeds times held-out batches pairs; (c) drop the $p$-value and
report the effect size with a bootstrap interval over seeds. Whichever is
chosen, write it into \S\,Evaluating in \texttt{paper.tex} before the
training runs are scheduled, because it sets the seed count.

\item \textbf{The token budget is inconsistent across arms.} \status{open}
\label{it:tokens}
\\
\emph{Problem.} The same quantity is set to four different values.
\begin{center}
\small
\begin{tabular}{@{}lll@{}}
\toprule
Where & Value & Applies to \\
\midrule
\texttt{paper.tex}, Model and training paragraph & 400 per turn & stated for all training \\
\texttt{paper.tex}, Appendix memory table & $T=400$ & memory arithmetic \\
\texttt{configs/loop\_grpo\_crystallizer.yaml} & 400 & arm C (single-turn GRPO) \\
\texttt{configs/loop\_grpo\_branched\_crystallizer.yaml} & 1000 per turn & arms E, F \\
\texttt{scripts/05\_train\_grpo\_branched.py} default & 400 & if the config omits it \\
Stage~1 and~2 as run & 600 & frozen arms A, B \\
Stage~1 and~2 re-run (planned) & 1000 & frozen arms A, B \\
\bottomrule
\end{tabular}
\end{center}
Stage~1 lost 28 of 96 controllers to the cut at 600 tokens and Stage~2
lost 44 of 96 probes. Arm~C at 400 would be handicapped against arm~A at
1000 and could fail the plumbing check for the wrong reason.
\\
\emph{Fix.} 1000 tokens per turn everywhere: the single-turn config, the
script default, the paper's paragraph, and the memory table (redo the
arithmetic at $T=1000$ or state that the table is illustrative at 400).
Arm~D's config, when it is written, inherits the same value.

\item \textbf{No probe-rate gate before Stages~4 to~6.} \status{open}
\label{it:probe-gate}
\\
\emph{Problem.} The probe is optional. If the frozen model under the
re-run protocol almost never probes, arms D, E and F collapse to arm~C and
the rollout tree carries no turn-1 signal to learn from. With $G_1=4$ and a
probe rate of ten percent, about a third of prompts have at least one
probing variant; at one percent, about four percent do. The extensions
section already names ten percent as the threshold for a prompt variant
to be trainable, but Stage~2's \emph{Decides} clause does not use it, and
it also does not say what happens if arm~B comes out below arm~A.
\\
\emph{Fix.} Add to Stage~2: a probe rate under the optional protocol of
at least ten percent is the gate for Stages~4 to~6; below it, the brief
is changed (not the training) and Stage~2 is re-run. Add the reading for
B below A: the probe run hurts, so the cause (budget, crashes, invented
results) is fixed in the brief before training.

\item \textbf{Checkpoint selection on the test set.} \status{open}
\label{it:checkpoint}
\\
\emph{Problem.} Results are ``taken at the best checkpoint'' and reported
on the 12 held-out batches. If the best checkpoint is chosen on those same
batches, the reported number is optimistically biased, and more so for arms
with noisier curves.
\\
\emph{Fix.} Choose the checkpoint on training reward or on a separate
validation draw of batches, then evaluate the held-out set once. Write the
rule into \S\,Evaluating.

\item \textbf{The evaluation set is thin.} \status{open}
\label{it:evalset}
\\
\emph{Problem.} Fast success depends strongly on where the limit sits in
$[0.08, 0.25]$: at the bottom the blind controller is already fast, at the
top the gap is largest. Twelve batches cover that range coarsely, and the
per-batch rate is itself a sampled quantity at temperature $0.8$.
\\
\emph{Fix.} 36 to 48 held-out batches with limits stratified across the
range (equal counts per tercile at least). Evaluation is inference only,
so this is cheap. Report fast success per tercile as well as overall.

\item \textbf{Add a frozen arm that is told the true limit.} \status{open}
\label{it:informed-arm}
\\
\emph{Problem.} If arm~F fails, nothing in the current design says whether
the model cannot write the ride-below-limit controller or cannot probe.
The Stage~2 record at 600 tokens already shows the controller itself
failing: setpoints of $0.18$ to $0.20$ with no margin, 24 of 96 batches
ruined after a median of 20 steps.
\\
\emph{Fix.} One inference-only arm: the frozen model, single turn, with
the true $\sigma_{\max}$ written into the brief. It is the model analogue
of Stage~0's informed side. Its fast-success rate is the ceiling any
probe-aware arm can reach with this model's controller-writing, and
subtracting it from the oracle gives the controller gap; subtracting arm~B
from it gives the probe gap. Add it to the ablation table as A$^{+}$ and
to Stage~1.

\item \textbf{Arms A to D form a two-by-two, not a chain.} \status{open}
\label{it:2x2}
\\
\emph{Problem.} The factors are frozen versus trained and one turn versus
two. The ablation table's ``what it isolates over the previous arm'' entry
for arm~C compares it to arm~B, which changes both factors at once.
\\
\emph{Fix.} Present A, B, C, D as a grid. C against A is the effect of
training without a probe; D against B is the effect of training with one;
B against A and D against C are the effect of the probe, frozen and
trained. E and F then extend D. Rewrite the table's last column
accordingly.

\item \textbf{Stages~4, 5 and~6 are one experiment, not three rungs.}
\status{open}
\label{it:one-experiment}
\\
\emph{Problem.} No stop condition is stated between Stage~4 and Stage~5
or between~5 and~6, and none makes sense: the question is the comparison,
and the arms share seeds. Only Stage~3 is a real gate, and it has no
numeric threshold, unlike Stages~0 and~0b.
\\
\emph{Fix.} Give Stage~3 a gate: arm~C reaches the blind reference (mean
reward about $+27$, fast success about $4$ to $7.5$ percent) within the
step budget on at least two of three seeds; failing that, the reward or the
plumbing is wrong and nothing else runs. Merge Stages~4 to~6 into one
stage with three arms run together on the same seeds, and keep the
\emph{Decides} text of the current Stage~6 as its decision rule.

\item \textbf{The stress set contradicts the brief.} \status{open}
\label{it:stress}
\\
\emph{Problem.} The brief tells the model the limit lies between $0.08$
and $0.25$. A stress set with limits outside that range therefore tests
what the model does when its stated prior is wrong, which is a different
question from robustness.
\\
\emph{Fix.} Either widen the range stated in the brief so the stress set
lies inside it, or keep the stress set and say in \S\,Evaluating that it
measures behaviour under a wrong prior. Either way, drop the word
``planned'' or schedule it.

\end{enumerate}

\section{Consistency fixes in the documents}
\label{sec:docs}

\begin{enumerate}[leftmargin=1.8em, itemsep=0.8em, resume]

\item \textbf{A stale mechanism check.} \status{open}
\label{it:mechanism}
\\
\emph{Problem.} The first diagnostic in \S\,Evaluating says that when
turn~1 emits a policy its reward should improve faster under~F. That
describes the old protocol, where a turn-1 policy was rehearsed and its
reward fed back. Under the current protocol a turn-1 controller ends the
task.
\\
\emph{Fix.} Replace it with: the probe rate over training, and the reward
of probed against unprobed trajectories. Under~F the probed trajectories'
turn-1 advantage should become positive before the probe rate rises;
under~D it cannot be separated.

\item \textbf{The tree weights a turn-1 controller as a full branch, and the
paper does not say so.} \status{open}
\label{it:replicate}
\\
\emph{Problem.} In \texttt{src/branched\_trainer.py} (the block commented
``The model wrote the controller in turn 1''), a first reply that is a
controller is entered into the reward grid $G_2$ times, so $\widehat V_0$
weights it like a probing variant and its turn-1 tokens appear $G_2$ times
in the loss. This is a reasonable choice, but it zeroes the within-prefix
variance for those variants, which inflates the between/within signal
ratio diagnostic whenever some variants probe and others do not.
\\
\emph{Fix.} State the choice in the Branched rollout subsection. Compute
the signal ratio over probing variants only, or report it alongside the
probe rate so the reader can tell the two apart.

\item \textbf{Matched leaves is not matched first-turn samples.}
\status{open}
\label{it:matched}
\\
\emph{Problem.} Arm~D draws 16 independent first turns per prompt; arms~E
and~F draw 4. That is exactly what E against D is meant to expose, but the
table caption says only that leaf counts match.
\\
\emph{Fix.} One sentence in the caption: D has $G$ independent first-turn
samples, E and F have $G_1$, and the E--D comparison is the price of that
reduction against the saving in generate calls.

\item \textbf{\texttt{methodology.tex} describes the old protocol.}
\status{open}
\label{it:methodology}
\\
\emph{Problem.} It has typed \texttt{policy}/\texttt{main} fences, a
rehearsal turn with the reward fed back, and an error floor of $-200$.
\texttt{paper.tex} has content-defined replies, an optional probe, and a
floor of $-300$. The two contradict each other.
\\
\emph{Fix.} Retire \texttt{methodology.tex}, or add a line under its title
saying it is superseded by \texttt{paper.tex} as of 2026-10-26 and is kept
for the derivation of the per-turn loss decomposition only.

\end{enumerate}

\section{Order of work}

\begin{enumerate}[leftmargin=1.8em]
  \item Items~\ref{it:tokens} and~\ref{it:probe-gate}, because the Stage~1
  and~2 re-runs need the 1000-token budget and the probe-rate gate is read
  off their result.
  \item Items~\ref{it:wilcoxon}, \ref{it:checkpoint} and~\ref{it:evalset},
  because they fix the seed count, the checkpoint rule and the evaluation
  set before any training run is scheduled.
  \item Item~\ref{it:informed-arm}, inference only, alongside the Stage~1
  re-run.
  \item Items~\ref{it:2x2}, \ref{it:one-experiment} and~\ref{it:stress},
  text changes to the plan, before Stage~3.
  \item Items~\ref{it:mechanism} to~\ref{it:methodology}, document
  consistency, at any time.
\end{enumerate}

\end{document}
